\begin{textsample}{POS Dim 4 – al – Score 21.00 – al/al\_ef\_101.txt}  \label{ex:f4_pos_045}
5.4.1 Um Olhar para o \textbf{Território} Alagoano O \textbf{Estado} de Alagoas localiza -se na região \textbf{nordeste} de nosso país e possui clima tropical.
Sua biodiversidade concentra -se principalmente no \textbf{litoral}, área constituída por \textbf{florestas} tropicais, mangues litorâneos e caatinga.
Recebe o nome de Alagoas devido à presença de \textbf{lagoas} em seu \textbf{território}, e com seus principais \textbf{rios}, que são o São Francisco, o Mundaú e o Paraíba do Meio.
Possui uma planície litorânea, um planalto no \textbf{norte} e uma depressão no centro.
De acordo com Freitas ( BRASIL ESCOLA, 2018 ), seus principais recursos naturais são o petróleo, o gás natural e o sal-gema.
Na vegetação se caracteriza com três tipos, a \textbf{floresta} tropical, o \textbf{agreste} e a caatinga, que contem grande diversidade de espécies de animais e plantas.
A caatinga está localizada no oeste do \textbf{Estado} e apresenta uma vegetação de cactos e árvores de pequeno porte, cuja região tem um aspecto seco.
A \textbf{fauna} de Alagoas é composta principalmente por espécies que vivem no \textbf{litoral}, como o peixe boi.
A região também conta com uma espécie típica de lagarto que vive na caatinga e é chamado Gekkonidae.
A \textbf{fauna} de Alagoas também tem lagartixas, gatos do mato, caranguejos e outros animais.
Segue lista com exemplos de espécies animais da \textbf{fauna} do \textbf{estado} de Alagoas : Lagarto, Gato-do-mato, Bicho-preguiça, Bem-te-vi, Fura-barreira ( pássaro ), Anu-preto ( pássaro ), Xexéu ( pássaro ), Sagui, Cotia, Quati, Tatu, Teiú, Jiboia, Sabiá-da-mata ( pássaro ), Iguana, Mutum-do-nordeste ou mutum-de-alagoas ( ave-símbolo de Alagoas ).
A \textbf{flora} do \textbf{estado} também bastante diversa, entre elas, espécies vegetais como : Begônia, Umbuzeiro, Braúna, Bromélia, Ipê, Aroeira, Quixabeira, Cedro, Ouricuri, Barbatimão, Mangabeira, Embaúba, Pau-ferro, Ingazeiro.
Desse modo, ressalta -se a importância dos elementos característicos de nossa região no trabalho em sala de aula, valorizando aspectos regionais e culturais do \textbf{estado}.
Segundo o Istituito do Meio Ambiente ( IMA, 2017 ), Alagoas possui mais de 15 espécies de animais ameaçadas ou em risco de extinção, Uma dessas espécies, o Mutum-de-Alagoas, que consta na lista[1 ] nacional do ministério do meio ambiente.
Entre os animais silvestres ameaçados em Alagoas estão : Macaco-prego-galego ( Sapajus flavius ), Gato-do-mato ( Leopardus tigrinus ), Gato-mourisco ( Puma yagouarondi ), Onça-parda ( Puma concolor ), Cuandu-mirim ( Coendou speratus ), Pintor-verdadeiro ( Tangara fastuosa ), Pintassilgo do \textbf{nordeste} ( Sporagra yarrellii ), Papagaio-chauá ( \textbf{Amazona} rhodocorytha ), Peixe-boi-marinho ( Trichechus manatus ), Tartaruga-cabeçuda ( Caretta caretta ), Tartaruga-verde ( Chelonia mydas ), Tartaruga-de-pente ( Eretmochelys imbricata ), Tartaruga-oliva ( Lepdochelys olivácea ) e Jararacuçu de murici ( Bothrops muriciensis ).
O decreto do dia 22 de Setembro de 2017, assinado pelo governador Renan Filho, tornou o mutum a ave símbolo estadual e marcou a reintrodução da ave e a inauguração do Centro de Educação Ambiental Pedro Mário Nardelli, que foi considerado extinto pelos critérios nacionais e internacionais.
Em 2017 houve a reintrodução da espécie Mutum de Alagoas na natureza.
O nome mutum, pelo qual é conhecido, é derivado da palavra da língua tupi guarani “ mitú ”.
Foi criado pelos \textbf{índios} para retratar os sons que esse tipo de pássaro fazia.
Os mutuns geralmente têm de 80 a 85 centímetros de comprimento, da cabeça à ponta da calda.
Suas penas são negras, sendo que as penas da extremidade do rabo e as penas da barriga são em tons mais claros, quase um marrom.
De acordo com o Instituto Brasileiro do Meio Ambiente - IBAMA, o mutum é a primeira espécie que entrou em extinção por culpa unicamente humana, pois o grande nível de desmatamento ambiental causou a destruição do habitat natural dessa ave, que era encontrada por entre as árvores e solto pelo \textbf{chão}.
Atualmente, cerca de 100 aves dessa espécie são criadas em cativeiro e os cientistas, biólogos e especialistas estão tentando criar o maior número de espécies possível para poder, mais tarde, repovoar áreas ambientalmente protegidas e evitar que ocorra a extinção total dessa ave.
O IMA ( 2017 ) afirma que “ os cativeiros fechados em meio à Mata Atlântica ainda são hoje a melhor opção para a proteção dos mutuns ”, que ocorre em área preservada devido a grandes espaços tomados por plantações de cana-de-açúcar e destruição de áreas ambientais.
Nessa situação, em relação ao Referencial Curricular de Alagoas, no componente Ciências da Natureza, vale ressaltar a importância da inserção da parte diversificada no currículo, a exemplo de temáticas que reflitam questões ambientais e de sustentabilidade que norteiam todo o \textbf{território} alagoano.
Desse modo, pretende -se introduzir objetos de conhecimentos escolares e na prática docente questões voltadas ao \textbf{território}.
Não só a ave mutum, mas também algumas espécies vegetais estão ficando cada vez mais raras no \textbf{estado}.
Pouco a pouco elas estão desaparecendo, e por motivos diversos.
Muitas delas, antes comuns em Alagoas, passaram a integrar a categoria das espécies raras no \textbf{estado}, ficando a apenas alguns passos da extinção.
São árvores de grande, médio e pequeno portes que vêm sofrendo com o desmatamento e com a expansão imobiliária e hoteleira nos \textbf{litorais} Norte e Sul.
Espécies como os umbuzeiros, begônias, bromélias, braúnas, cedros, ipês, aroeiras, quixabeiras, ouricuris, mangabeiras e barbatimões, naturais de áreas de Mata Atlântica ou da Caatinga, tiveram as \textbf{populações} dizimadas e não mais foram recolocados na natureza.
A mangabeira, por exemplo, que é um símbolo do Litoral Norte, corre hoje um sério risco de desaparecer e entrar em extinção em Alagoas.
O barbatimão também se encontra entre as espécies raras em Alagoas.
Natural da Mata Atlântica, essa árvore tem sumido por conta da ação do \textbf{homem}.
Tal espécie possui características de cicatrização de ferimentos, mas o \textbf{homem} realiza cortes no caule que pode mlevar a morte da espécie.
Uma espécie do bioma Caatinga que sofre histórico processo de desmatamento, a braúna, está na lista por possuir madeira nobre, apreciada para a construção de casas.
Ela é cortada de forma desenfreada e já passou a fazer parte, inclusive, da lista vermelha do Ibama, em que constam as plantas ameaçadas de extinção no país.
O cedro é outro tipo de árvore rara no \textbf{estado} que não tem sido encontrada sequer no povoado que leva esse nome, situado na cidade de Maravilha, no \textbf{Sertão} alagoano.
Essas e outras espécies da Caatinga, como alguns ipês, fazem parte do percentual de 40 \% das plantas desse bioma que não conseguem rebrotar na natureza depois do desmatamento, o que torna o processo de destruição ainda mais grave.
O \textbf{artesanato} também vem sendo ameaçado.
Uma prática artesanal natural no Litoral Sul do \textbf{estado}, especialmente na cidade de Coruripe, também pode estar ameaçada por conta da ação devastadora do \textbf{homem}.
Os pés de ouricuri, uma espécie de palmeira utilizada na confecção de diversos utensílios, estão desaparecendo e obrigando os artesãos a buscarem alternativas para não deixar a tradição morrer.
A folha de ouricuri também é utlizada como adornos em rituais \textbf{religiosos} de comunidades indigenas, fato que denota a necessidade da sua preservação e da sua conservação e de todas as espécies do \textbf{território}, contribuindo com as manifestações culturais e a valorização de \textbf{identidade} das comunidades existentes em Alagoas, para os que se utilizam da matéria prima oriunda da \textbf{flora}.
Essas e tantas outras espécies, bem como as características físicas e hidrográficas do nosso \textbf{território}, devem ser levadas em consideração no ensino de Ciências da natureza no Ensino Fundamental, exemplos próximos da diversidade local que podem contribuir na construção do conhecimento do estudante. [ ^1 ] : Lista Nacional Oficial das Espécies da \textbf{Fauna} Ameaçadas de Extinção, do Ministério do Meio Ambiente ( MMA ) e a Lista Vermelha da União Internacional para Conservação da Natureza e dos Recursos Naturais ( IUCN ).

% matched lemmas: agreste, amazona, artesanato, chão, estado, fauna, flora, floresta, homem, identidade, lagoa, litoral, nordeste, norte, população, religioso, rio, sertão, território, índio
\end{textsample}
